% ============================================================================
% Fichier  : partie3/P03_C09_reseau.tex
% Titre    : Attaques Réseau et Authentification
% Auteur   : MINKA MI NGUIDJOÏ Thierry Emmanuel
% Version  : 1.0 -- novembre 2026
% ============================================================================
% Sources : RFC 4795 (LLMNR) ; RFC 4120 (Kerberos) ; documentation publique
% de Microsoft sur NTLM et la signature SMB ; MITRE ATT&CK (T1557). Résultats
% de l'Opération LOGONE : cas fictif.
% ============================================================================

\chapter{Attaques Réseau et Authentification}
\label{chap:reseau}
\chaptermark{Attaques réseau et authentification}

\epigraph{%
    « Un réseau qui crie ses questions à tout le monde\\
    finit par recevoir une mauvaise réponse. »%
}{Principe du cours}

\niveauChapitre{Intermédiaire à expert}{Chapitre~\ref{chap:exploitation}~;
cours SEC3266, chapitres sur les réseaux et les identités.}

\begin{boxobjectifs}
\begin{itemize}
    \item Expliquer les protocoles d'authentification d'un réseau Windows
          (NTLM, Kerberos).
    \item Comprendre l'empoisonnement des résolutions de noms locales.
    \item Comprendre le relais d'authentification et ses conditions.
    \item Relier chaque attaque à sa détection et à sa parade.
    \item Rédiger les constats correspondants.
\end{itemize}
\end{boxobjectifs}

\bigskip

% ============================================================================
\section*{Question directrice : comment un poste ordinaire peut-il recueillir des secrets sans rien attaquer~?}
% ============================================================================

Le 26~octobre, depuis le poste fourni, Mme~NGO~MBOCK observe le trafic
local de Garoua. Des postes demandent à tout le réseau « qui est
srv-imprim~? », un serveur qui n'existe plus. Une machine qui répondrait
« c'est moi » recevrait leur tentative d'authentification. Ce chapitre
explique pourquoi, et comment fermer cette porte.

\medskip
\textbf{Plan du chapitre.}
\begin{enumerate}
    \item NTLM et Kerberos (section~\ref{sec:c09-protocoles}).
    \item L'empoisonnement des résolutions locales (section~\ref{sec:c09-llmnr}).
    \item Le relais d'authentification (section~\ref{sec:c09-relais}).
    \item Détecter et corriger (section~\ref{sec:c09-parades}).
\end{enumerate}

\separateur

% ============================================================================
\section{NTLM et Kerberos}
\label{sec:c09-protocoles}
% ============================================================================

\begin{tableaucours}{L{2.6cm}L{4.8cm}L{4.8cm}}{%
    \tableauHeader{Protocole} &
    \tableauHeader{Principe} &
    \tableauHeader{Faiblesse exploitable}
}
NTLM & Défi et réponse calculée à partir du haché du mot de passe & La
réponse peut être relayée ou cassée hors ligne si le mot de passe est
faible \\
Kerberos & Tickets délivrés par un centre de distribution de clés
(RFC~4120) & Tickets de service chiffrés avec le secret du compte de
service (Chapitre~\ref{chap:ad}) \\
\end{tableaucours}
\vspace{-0.8em}
\begin{center}{\small\itshape Tableau~9.1. Deux protocoles d'authentification}\end{center}

Kerberos, défini par la RFC~4120, est le protocole par défaut d'un domaine
Active Directory~\cite{rfc4120}~; NTLM subsiste pour la compatibilité, et
c'est souvent lui que vise l'attaquant.

% ============================================================================
\section{L'Empoisonnement des Résolutions Locales}
\label{sec:c09-llmnr}
% ============================================================================

Quand un nom ne se résout pas par le DNS, Windows peut interroger tout le
réseau local par LLMNR (RFC~4795) ou NBT-NS~\cite{rfc4795}. N'importe quelle
machine peut répondre. Un attaquant présent sur le réseau répond à ces
questions et reçoit en retour une tentative d'authentification NTLM.

\begin{boxdef}
\termecle{Empoisonnement de résolution}~: réponse frauduleuse à une
requête de résolution de nom diffusée sur le réseau local, qui détourne la
connexion de la victime vers l'attaquant (ATT\&CK, T1557).
\end{boxdef}

\begin{boxscenario}{Opération LOGONE. Les questions sans réponse}
\textbf{Constatation.} En une heure d'écoute passive, 37~requêtes LLMNR
pour des noms inexistants (dont « srv-imprim », retiré en 2025), émises par
22~postes. Avec l'accord de M.~KOUNOU, une réponse contrôlée de dix minutes
recueille 6~réponses NTLM de 4~comptes~; aucune n'est cassée, les mots de
passe concernés étant longs.

\medskip
\textbf{Analyse} (confiance élevée). La faille est structurelle (LLMNR et
NBT-NS actifs)~; son impact dépend de la force des mots de passe et de la
possibilité de relais.
\end{boxscenario}

% ============================================================================
\section{Le Relais d'Authentification}
\label{sec:c09-relais}
% ============================================================================

Au lieu de casser la réponse NTLM, l'attaquant peut la \emph{relayer}
aussitôt vers un autre serveur, et s'y authentifier comme la victime. Le
relais est possible lorsque le serveur cible n'exige pas la
\termecle{signature} des échanges (SMB) ou la liaison au canal (HTTP,
LDAP).

\begin{tableaucours}{L{4.0cm}L{8.2cm}}{%
    \tableauHeader{Condition} &
    \tableauHeader{Situation chez LiZbiZ}
}
Requêtes diffusées détournables & LLMNR et NBT-NS actifs \\
Serveurs sans signature SMB exigée & 31 serveurs sur 34 à Garoua \\
Comptes à privilèges qui s'authentifient & Comptes d'administration
utilisés depuis des postes ordinaires \\
\end{tableaucours}
\vspace{-0.8em}
\begin{center}{\small\itshape Tableau~9.2. Les conditions du relais (Opération LOGONE)}\end{center}

\begin{boxalert}
\textbf{Violation typique.} Relayer une authentification vers un serveur de
production sans accord.

\smallskip
\textbf{Règle~:} le relais crée une session réelle au nom d'un salarié~:
il n'est démontré que vers une cible et un compte convenus, et la session
est fermée aussitôt la preuve faite.
\end{boxalert}

% ============================================================================
\section{Détecter et Corriger}
\label{sec:c09-parades}
% ============================================================================

\begin{tableaucours}{L{4.4cm}L{7.8cm}}{%
    \tableauHeader{Parade} &
    \tableauHeader{Effet}
}
Désactiver LLMNR et NBT-NS & Supprime les requêtes diffusées détournables \\
Exiger la signature SMB & Empêche le relais vers SMB \\
Protection étendue de l'authentification (LDAP, HTTP) & Lie
l'authentification au canal \\
Réduire NTLM au profit de Kerberos & Réduit la surface \\
Comptes d'administration séparés, bastion & Les secrets d'administration ne
circulent pas depuis les postes \\
Détection & Répondant inconnu aux requêtes LLMNR~; authentifications NTLM
d'un compte depuis une machine inhabituelle \\
\end{tableaucours}
\vspace{-0.8em}
\begin{center}{\small\itshape Tableau~9.3. Parades et détection}\end{center}

\begin{boxobjet}{T-07 · Signature SMB non exigée et LLMNR actif}
\textbf{Description~:} 31~serveurs sur 34 n'exigent pas la signature SMB~;
LLMNR et NBT-NS sont actifs sur les postes. \textbf{Preuve~:} relevés du
26~octobre~; réponses NTLM recueillies (masquées). \textbf{Gravité~:}
moyenne (élevée si un compte à privilèges est capturé).
\textbf{Recommandation~:} désactiver LLMNR et NBT-NS par stratégie de
groupe~; exiger la signature SMB. \textbf{Détection~:} surveiller les
réponses LLMNR d'hôtes non autorisés.
\end{boxobjet}

\begin{boxmetier}{Analyste SOC et équipe rouge}
\textbf{Poser un piège.} Un poste « leurre » émet régulièrement une requête
LLMNR pour un nom inexistant~: toute réponse signale un empoisonneur sur le
réseau.
\end{boxmetier}

\separateur

\begin{boxresume}
\textbf{Ce qu'il faut retenir de ce chapitre~:}
\begin{itemize}
    \item \textbf{NTLM} se relaie ou se casse si le mot de passe est faible~;
          \textbf{Kerberos} est le protocole par défaut du domaine.
    \item \textbf{LLMNR et NBT-NS} diffusent des questions auxquelles
          n'importe qui peut répondre.
    \item \textbf{Le relais} exige l'absence de signature ou de liaison au
          canal.
    \item \textbf{Désactiver, signer, lier, séparer, détecter}~: les parades.
    \item \textbf{Démontrer} seulement vers une cible convenue, et fermer
          aussitôt.
\end{itemize}
\end{boxresume}

\begin{boxinfo}
\textbf{Points de contrôle de l'élève.} Je sais~: comparer NTLM et
Kerberos~; expliquer l'empoisonnement LLMNR~; énoncer les conditions du
relais~; proposer les parades et une détection.
\end{boxinfo}

\bigskip

\begin{boxexo}[title={\ding{45}\quad Exercice~9.1. LLMNR \niveaubase}]
Pourquoi un poste envoie-t-il une requête LLMNR pour « srv-imprim »~?
Quelle mesure d'hygiène l'aurait évité~?
\end{boxexo}

\begin{boxexo}[title={\ding{45}\quad Exercice~9.2. Proportion \niveaubase}]
Quelle part des serveurs de Garoua n'exige pas la signature SMB~?
\end{boxexo}

\begin{boxexo}[title={\ding{45}\quad Exercice~9.3. Kerberos \niveaubase}]
Quel organisme a défini Kerberos version~5 dans la RFC~4120~? Quel est le
rôle du centre de distribution de clés~?
\end{boxexo}

\begin{boxexo}[title={\ding{45}\quad Exercice~9.4. Casser ou relayer \niveaumoyen}]
Pourquoi les réponses recueillies n'ont-elles pas été cassées~? Que
faudrait-il pour qu'elles le soient~?
\end{boxexo}

\begin{boxexo}[title={\ding{45}\quad Exercice~9.5. Les conditions \niveaumoyen}]
Laquelle des trois conditions du Tableau~9.2 est la plus simple à
supprimer~? Laquelle protège le mieux~?
\end{boxexo}

\begin{boxexo}[title={\ding{45}\quad Exercice~9.6. Le piège \niveaumoyen}]
Décrivez le poste « leurre » et la règle d'alerte du SOC.
\end{boxexo}

\begin{boxexo}[title={\ding{45}\quad Exercice~9.7. Effets de bord \niveaumoyen}]
Quels risques de dysfonctionnement faut-il vérifier avant d'exiger la
signature SMB partout~?
\end{boxexo}

\begin{boxexo}[title={\ding{45}\quad Exercice~9.8. ATT\&CK \niveaumoyen}]
Dans quelle tactique d'ATT\&CK se range l'empoisonnement de résolution~?
Pourquoi~?
\end{boxexo}

\begin{boxexo}[title={\ding{45}\quad Exercice~9.9. Plan de correction \niveauexpert}]
Rédigez le plan de correction de T-07 en trois vagues, avec les tests de
non-régression.
\end{boxexo}

\begin{boxexo}[title={\ding{45}\quad Exercice~9.10. Réduire NTLM \niveauexpert}]
Proposez une démarche pour mesurer l'usage de NTLM chez LiZbiZ et le
réduire sans interrompre les applications.
\end{boxexo}

\bigskip

\begin{boxinfo}
\textbf{Pour aller plus loin.}
\begin{itemize}
    \item RFC~4120 (Kerberos)~\cite{rfc4120}~; RFC~4795 (LLMNR)~\cite{rfc4795}.
    \item MITRE ATT\&CK~\cite{mitreattack}.
\end{itemize}
\end{boxinfo}

\begin{boxinfo}
\textbf{Transition.} Kerberos et l'annuaire forment le cœur du réseau
Windows~: le Chapitre~\ref{chap:ad} est consacré à Active Directory.
\end{boxinfo}
